% AM I 09, 18; AM II 02, 04, 07, 09–10.
\chapter{Szeregi potęgowe}
\label{ma2:ch:power-series}

Niech $\K\in\{\R,\C\}$, $E$ będzie przestrzenią
Banacha nad $\K$, a $z_0\in\K$. Rozważamy szeregi
o współczynnikach wektorowych i zmiennej skalarnej.
Przypadek $E=\K$ daje zwykłe szeregi potęgowe.

\section{Promień zbieżności}

\begin{definicja}
\label{ma2:def:power-series}
Szeregiem potęgowym o środku $z_0$ nazywamy
szereg funkcyjny
\[
 \sum_{n=0}^\infty a_n(z-z_0)^n,\qquad a_n\in E.
\]
Przyjmujemy $w^0=1$ również dla $w=0$.
Liczbę
\[
 R=\left(\limsup_{n\to\infty}\|a_n\|^{1/n}\right)^{-1}
 \in[0,+\infty]
\]
nazywamy promieniem zbieżności; przyjmujemy
$1/0=+\infty$ i $1/(+\infty)=0$.
\end{definicja}

\begin{twierdzenie}[Wzór Cauchy'ego--Hadamarda]
\label{ma2:thm:cauchy-hadamard}
Szereg potęgowy jest bezwzględnie zbieżny
dla $|z-z_0|<R$, a rozbieżny dla $|z-z_0|>R$.
W środku $z_0$ ma zawsze sumę $a_0$.
Dla każdego skończonego $0\le r<R$ szereg
jest zbieżny normalnie na zbiorze
$|z-z_0|\le r$. W szczególności jest
jednostajnie zbieżny na każdym zwartym
podzbiorze zbioru $|z-z_0|<R$.
\end{twierdzenie}
\begin{proof}
Dla $z=z_0$ wszystkie wyrazy poza $a_0$
są zerowe. Dla $\rho=|z-z_0|>0$ mamy
\[
 \limsup_n \|a_n(z-z_0)^n\|^{1/n}
 =\rho\limsup_n\|a_n\|^{1/n}.
\]
Mnożenie przez dodatnią skończoną liczbę
zachowuje granicę górną, także nieskończoną.
Jeśli $0<\rho<R$, prawa strona jest mniejsza
niż $1$, także gdy $R=+\infty$. Jeśli
$\rho>R$, jest większa niż $1$, także gdy $R=0$.
Stosujemy kryterium pierwiastkowe
(twierdzenie \ref{ma2:thm:vector-ratio-root};
wersja skalarna: AM I, rozdział 09,
etykieta \texttt{ma1:thm:root-test}).
Dla $0<r<R$ szereg liczbowy $\sum_n\|a_n\|r^n$
jest zbieżny z tego samego kryterium.
Na domkniętym kole lub przedziale promienia $r$
supremum normy $n$-tego wyrazu wynosi
$\|a_n\|r^n$, gdyż zbiór zawiera punkt $z_0+r$.
Dla $r=0$ supremum jest zerowe dla $n\ge1$
i równe $\|a_0\|$ dla $n=0$.
Kryterium Weierstrassa
\ref{ma2:thm:weierstrass-test} daje zbieżność
jednostajną. Jeśli $K$ jest niepustym zwartym
podzbiorem otwartego koła lub przedziału,
funkcja $z\mapsto|z-z_0|$ osiąga na $K$
maksimum $r<R$. Stosujemy wynik na zbiorze
promienia $r$ i ograniczamy do $K$.
Dla pustego $K$ warunek jest automatyczny.
\end{proof}

Zatem dla $0<R<+\infty$ zbieżność zachodzi
wewnątrz koła (przedziału w przypadku rzeczywistym),
a na brzegu wymaga osobnego badania.
Promień $R$ jest również supremum zbioru
liczb $r\ge0$, dla których $\sum_n\|a_n\|r^n$
jest skończony. Wynika to z obu części twierdzenia,
z uwzględnieniem konwencji w zerze.

\begin{wniosek}[Oszacowanie reszty]
\label{ma2:cor:power-remainder}
Jeśli $0\le r<s<R$ i $s<\infty$, to przy
$C_s=\sum_{n\ge0}\|a_n\|s^n$ zachodzi
\[
 \sup_{|z-z_0|\le r}
 \left\|f(z)-\sum_{n=0}^Na_n(z-z_0)^n\right\|
 \le C_s(r/s)^{N+1},\qquad N\ge0,
\]
gdzie $f$ jest sumą szeregu.
\end{wniosek}
\begin{proof}
Dla $n>N$ mamy
$\|a_n\|r^n=\|a_n\|s^n(r/s)^n
\le\|a_n\|s^n(r/s)^{N+1}$.
Sumujemy nierówności i stosujemy oszacowanie
reszty z twierdzenia \ref{ma2:thm:weierstrass-test}.
\end{proof}

\begin{wniosek}[Kryterium ilorazowe]
\label{ma2:cor:power-ratio}
Jeżeli $a_n\ne0$ od pewnego indeksu oraz
$\lim_n\|a_{n+1}\|/\|a_n\|=L\in[0,+\infty]$,
to $R=1/L$ z powyższymi konwencjami.
\end{wniosek}
\begin{proof}
Dla $\rho>0$ iloraz norm kolejnych wyrazów
szeregu w punkcie $z_0+\rho$ ma granicę $\rho L$.
Z twierdzenia \ref{ma2:thm:vector-ratio-root}
szereg zbiega bezwzględnie dla $\rho L<1$,
a jest rozbieżny dla $\rho L>1$.
Porównując z twierdzeniem
\ref{ma2:thm:cauchy-hadamard}, otrzymujemy
$R=1/L$, ponieważ różne promienie wyznaczałyby
sprzeczne wnioski dla pewnego dodatniego $\rho$
ściśle pomiędzy nimi. Gdy $L=0$, wszystkie
dodatnie $\rho$ dają zbieżność, a gdy $L=+\infty$,
żadne dodatnie $\rho$ nie daje zbieżności.
\end{proof}

\begin{przyklad}
\label{ma2:ex:power-radii}
Szeregi $\sum_{n\ge0}n!z^n$, $\sum_{n\ge0}z^n$
oraz $\sum_{n\ge0}z^n/n!$ mają odpowiednio
promienie $0$, $1$ i $+\infty$. Ilorazy norm
współczynników wynoszą odpowiednio $n+1$,
$1$ oraz $1/(n+1)$, więc stosujemy wniosek
\ref{ma2:cor:power-ratio}.
W trzecim przypadku otrzymujemy funkcję
wykładniczą z AM I, rozdział 18, definicja
18.1.1; jej definicji nie wprowadzamy ponownie.
Wielomian ma promień nieskończony, gdyż
ciąg jego współczynników jest od pewnego
indeksu zerowy.
\end{przyklad}

\section{Zbieżność na brzegu}

\begin{przyklad}[Trzy zachowania na brzegu]
\label{ma2:ex:power-boundary}
Każdy z szeregów
\[
 \sum_{n\ge1}z^n,\qquad
 \sum_{n\ge1}\frac{z^n}{n},\qquad
 \sum_{n\ge1}\frac{z^n}{n^2}
\]
ma promień $1$ z wniosku \ref{ma2:cor:power-ratio}.
Pierwszy jest rozbieżny we wszystkich punktach
$|z|=1$, ponieważ normy wyrazów są równe $1$.
Trzeci jest zbieżny bezwzględnie w każdym
takim punkcie i normalnie na całym domkniętym
kole, ponieważ $\sum 1/n^2<\infty$.
Drugi jest rozbieżny w $z=1$ (szereg harmoniczny).
Dla $|z|=1$, $z\ne1$, zachodzi
\[
 \left|\sum_{n=1}^k z^n\right|
 =\left|\frac{z(1-z^k)}{1-z}\right|
 \le\frac{2}{|1-z|}.
\]
Z kryterium Dirichleta
\ref{ma2:thm:uniform-dirichlet} na zbiorze
jednoelementowym wynika zbieżność. Nie jest
ona bezwzględna, gdyż $\sum |z^n/n|=\sum1/n$.
Na każdym podzbiorze okręgu oddalonym
od $1$ o dodatnią odległość stała w powyższym
oszacowaniu jest wspólna, więc zbieżność
jest jednostajna. W przypadku rzeczywistym
drugi szereg zbiega warunkowo w $-1$
i jest rozbieżny w $1$.
\end{przyklad}

\begin{figure}[htbp]
\centering
\begin{tikzpicture}[x=1.5cm,y=1.5cm,>=Stealth]
 \foreach \p/\opis in {0/{A},3/{B},6/{C}}{
  \begin{scope}[shift={(\p,0)}]
   \fill[blue!10] (0,0) circle[radius=1];
   \draw[->,gray] (-1.15,0)--(1.2,0);
   \draw[->,gray] (0,-1.15)--(0,1.2);
   \node[below] at (0,-1.2) {\opis};
   \node[font=\scriptsize,above left] at (0,0) {$0$};
  \end{scope}
 }
 \draw[red!75!black,very thick,dashed] (0,0) circle[radius=1];
 \draw[orange!85!black,very thick] (3,0) circle[radius=1];
 \filldraw[fill=white,draw=red!75!black,thick] (4,0) circle[radius=2pt];
 \draw[blue!70!black,very thick] (6,0) circle[radius=1];
\end{tikzpicture}

\small
A: $\sum z^n$; \quad
B: $\sum z^n/n$; \quad
C: $\sum z^n/n^2$.
\caption{Koła o promieniu zbieżności $1$.
Wewnątrz zachodzi zbieżność bezwzględna.
Brzeg przerywany oznacza rozbieżność,
pomarańczowy --- zbieżność warunkową
(z wyłączeniem punktu $1$), niebieski ---
zbieżność bezwzględną.}
\label{ma2:fig:power-boundary}
\end{figure}
\FloatBarrier

\begin{twierdzenie}[Twierdzenie Abela o granicy radialnej]
\label{ma2:thm:power-abel-boundary}
Niech $0<R<+\infty$ i $|\zeta-z_0|=R$.
Jeżeli szereg $\sum_n a_n(\zeta-z_0)^n$
jest zbieżny do $S$, to
\[
 \lim_{t\to1^-}f(z_0+t(\zeta-z_0))=S.
\]
Co więcej, szereg po podstawieniu
$z=z_0+t(\zeta-z_0)$ zbiega jednostajnie
dla $t\in[0,1]$.
\end{twierdzenie}
\begin{proof}
Połóżmy $b_n=a_n(\zeta-z_0)^n$.
Szereg $\sum b_n$ rozpatrywany jako szereg
funkcji stałych jest jednostajnie zbieżny
na $[0,1]$. Ciąg mnożników $t^n$ jest
dla każdego $t\in[0,1]$ nierosnący
i ograniczony przez $1$.
Jednostajne kryterium Abela
\ref{ma2:thm:uniform-abel}, po przesunięciu
indeksu początkowego, daje jednostajną
zbieżność szeregu $\sum b_nt^n$ na $[0,1]$.
Jego suma jest ciągła z twierdzenia
\ref{ma2:thm:uniform-continuous-limit},
gdyż sumy częściowe są ciągłe.
W $t=1$ wynosi $S$, a w $0\le t<1$
jest równa $f(z_0+t(\zeta-z_0))$.
\end{proof}
Odwrotność twierdzenia nie zachodzi.
Dla $f(z)=1/(1+z)=\sum_{n\ge0}(-z)^n$
wewnątrz koła jednostkowego granica radialna
w $z=1$ wynosi $1/2$, lecz szereg na brzegu
$\sum_{n\ge0}(-1)^n$ jest rozbieżny.

\section{Różniczkowanie i całkowanie}

\begin{twierdzenie}[Promień szeregu pochodnych]
\label{ma2:thm:power-derivative-radius}
Szereg
$\sum_{n\ge1}na_n(z-z_0)^{n-1}$
ma ten sam promień zbieżności co
$\sum_{n\ge0}a_n(z-z_0)^n$.
\end{twierdzenie}
\begin{proof}
Niech $R'$ będzie promieniem szeregu pochodnych.
Jeżeli $0<r<s<R$, to zbiega szereg
$\sum_n\|a_n\|s^n$. Liczby
\[
 n r^{n-1}/s^n=\frac{n}{s}(r/s)^{n-1}
\]
są ograniczone, gdyż iloraz kolejnych
dodatnich liczb dąży do $r/s<1$,
więc od pewnego indeksu maleją.
Zatem $\sum_{n\ge1}n\|a_n\|r^{n-1}$
jest zbieżny. Dla każdego $r<R$ dodatniego
można wybrać takie $s$, więc $R'\ge R$
(również dla $R=0$ lub $R=+\infty$).
Odwrotnie, dla $0<r<R'$ mamy
\[
 \sum_{n\ge1}\|a_n\|r^n
 \le r\sum_{n\ge1}n\|a_n\|r^{n-1}<\infty.
\]
Po dodaniu $\|a_0\|$ wynika $R\ge R'$.
Gdy $R'=0$, ta nierówność jest automatyczna.
\end{proof}

\begin{twierdzenie}[Różniczkowanie wyraz po wyrazie]
\label{ma2:thm:power-derivative}
Jeżeli $R>0$, suma $f$ szeregu potęgowego
jest różniczkowalna w każdym punkcie
$|z-z_0|<R$ (w przypadku zespolonym
w sensie granicy zespolonych ilorazów różnicowych)
oraz
\[
 f'(z)=\sum_{n\ge1}na_n(z-z_0)^{n-1}.
\]
Dla każdego $k\ge1$ istnieje pochodna $k$-tego rzędu
i
\[
 f^{(k)}(z)=
 \sum_{n=k}^\infty\frac{n!}{(n-k)!}a_n(z-z_0)^{n-k}.
\]
Wszystkie te szeregi mają promień $R$.
\end{twierdzenie}
\begin{proof}
Ustalmy $z$ wewnątrz obszaru zbieżności.
Wybieramy skończone $r$ z $|z-z_0|<r<R$.
Dla $h\ne0$ dostatecznie małego mamy
$|z+h-z_0|\le r$. Dla $w=z-z_0$ i $n\ge1$
tożsamość różnicy potęg daje
\[
 \frac{(w+h)^n-w^n}{h}
 =\sum_{j=0}^{n-1}(w+h)^{n-1-j}w^j.
\]
Każdy składnik ma moduł nie większy niż
$r^{n-1}$, a dla ustalonego $n$ prawa
strona dąży do $nw^{n-1}$ przy $h\to0$.
Z twierdzenia \ref{ma2:thm:power-derivative-radius}
szereg $\sum_{n\ge1}n\|a_n\|r^{n-1}$ jest
zbieżny. Kryterium Weierstrassa
\ref{ma2:thm:weierstrass-test} zastosowane
do ilorazów różnicowych (jako funkcji $h$)
zapewnia jednostajną zbieżność ich szeregu
na dostatecznie małym przebitym otoczeniu zera.
Z liniowości granic sum częściowych suma
tego szeregu jest równa $(f(z+h)-f(z))/h$.
Z twierdzenia \ref{ma2:thm:interchange-limits},
stosowanego do sum częściowych, możemy
przejść do granicy przy $h\to0$ wyraz po wyrazie.
Otrzymujemy wzór na $f'$.
Kolejne wzory wynikają indukcyjnie:
szereg pochodnej ma ten sam promień
z twierdzenia \ref{ma2:thm:power-derivative-radius},
więc na każdym kroku stosujemy już
udowodnioną część twierdzenia. Współczynniki
mnożą się przez kolejny wykładnik potęgi,
co daje $n!/(n-k)!$.
\end{proof}

\begin{wniosek}[Jednoznaczność współczynników]
\label{ma2:cor:power-coefficients}
Dla szeregu o dodatnim promieniu zbieżności
zachodzi $a_k=f^{(k)}(z_0)/k!$ dla $k\ge0$,
gdzie $f^{(0)}=f$.
W szczególności dwa szeregi o tym samym środku
i dodatnich promieniach, których sumy zgadzają
się w otoczeniu środka, mają identyczne współczynniki.
\end{wniosek}
\begin{proof}
Wstawiamy $z=z_0$ do wzoru z twierdzenia
\ref{ma2:thm:power-derivative}; wszystkie składniki
poza $n=k$ znikają. Dla $k=0$ korzystamy
z $f(z_0)=a_0$. Równość funkcji w otoczeniu
środka pociąga równość wszystkich ich pochodnych
w środku, z definicji pochodnej i indukcji.
\end{proof}

\begin{wniosek}[Funkcja pierwotna i całkowanie]
\label{ma2:cor:power-primitive}
Jeśli $R>0$ i $c\in E$, szereg
\[
 F(z)=c+\sum_{n\ge0}\frac{a_n}{n+1}(z-z_0)^{n+1}
\]
ma promień $R$, a $F'=f$ dla $|z-z_0|<R$.
Dla rzeczywistej zmiennej i dowolnych
$u<v$ z $[u,v]\subseteq(z_0-R,z_0+R)$
zachodzi
\[
 \int_u^v f(t)\,\mathrm dt
 =\sum_{n\ge0}\frac{a_n}{n+1}
      ((v-z_0)^{n+1}-(u-z_0)^{n+1}).
\]
\end{wniosek}
\begin{proof}
Formalny szereg pochodnych szeregu $F$ jest
szeregiem $f$, więc równość promieni wynika
z twierdzenia \ref{ma2:thm:power-derivative-radius},
a wzór na pochodną z twierdzenia
\ref{ma2:thm:power-derivative}.
Na $[u,v]$ szereg $f$ zbiega jednostajnie.
Wniosek \ref{ma2:cor:termwise-integral} pozwala
całkować wyraz po wyrazie; całka wektora stałego
pomnożonego przez skalarną funkcję jest iloczynem
tego wektora i całki skalarnej, co wynika wprost
z sum Riemanna. Całkę potęgi obliczamy
z twierdzenia Newtona--Leibniza (AM I, 29.6.3).
\end{proof}

\section{Działania algebraiczne}

\begin{twierdzenie}
\label{ma2:thm:power-algebra}
Niech $f=\sum a_n(z-z_0)^n$ i
$g=\sum b_n(z-z_0)^n$ będą szeregami
o promieniach $R_f,R_g$, ze współczynnikami
odpowiednio w przestrzeniach Banacha $E,F$
nad tym samym ciałem $\K$.
Jeśli $T:E\times F\to G$ jest ciągłym
odwzorowaniem dwuliniowym do przestrzeni
Banacha $G$, to szereg
\[
 \sum_{n\ge0}c_n(z-z_0)^n,\qquad
 c_n=\sum_{j=0}^n T(a_j,b_{n-j}),
\]
ma promień co najmniej $\min(R_f,R_g)$
i wewnątrz tego obszaru sumuje się do $T(f(z),g(z))$.
Jeśli $E=F$, analogicznie dla $\alpha,\beta\in\K$
szereg o współczynnikach $\alpha a_n+\beta b_n$
ma promień co najmniej $\min(R_f,R_g)$
i sumę $\alpha f+\beta g$ w tym obszarze.
\end{twierdzenie}
\begin{proof}
Dla $0<r<\min(R_f,R_g)$ oba szeregi
$\sum\|a_n\|r^n$ i $\sum\|b_n\|r^n$ są zbieżne.
Z lematu \ref{ma2:lem:bilinear-bound}
istnieje $K\ge0$ z
$\|T(a,b)\|\le K\|a\|\|b\|$.
Dla każdego $N$ mamy
\[
 \sum_{n=0}^N\|c_n\|r^n
 \le K\left(\sum_{j\ge0}\|a_j\|r^j\right)
      \left(\sum_{k\ge0}\|b_k\|r^k\right).
\]
Wynika to z nierówności trójkąta,
$r^n=r^jr^{n-j}$ i rozszerzenia sumowania
na wszystkie pary nieujemnych indeksów.
Sumy częściowe szeregu norm są niemalejące
i ograniczone, zatem szereg norm jest zbieżny.
Wynika stąd dolne ograniczenie promienia;
gdy minimum jest zerem, jest ono automatyczne.
Dla $|z-z_0|<\min(R_f,R_g)$ stosujemy
wniosek \ref{ma2:cor:bilinear-absolute}
do obu szeregów w tym punkcie.
Dwuliniowość daje
$T(a_j(z-z_0)^j,b_k(z-z_0)^k)
=T(a_j,b_k)(z-z_0)^{j+k}$.
Zatem iloczyn Cauchy'ego jest dokładnie
wskazanym szeregiem. Dla kombinacji liniowej
stosujemy nierówność trójkąta do współczynników
oraz liniowość granic sum częściowych.
\end{proof}

\begin{przyklad}
\label{ma2:ex:power-cancellation}
Promień wynikowego szeregu może wzrosnąć.
Dla $f(z)=\sum_{n\ge0}z^n$ i $g(z)=1-z$
iloczyn Cauchy'ego ma współczynnik $1$ dla
$n=0$ i $1-1=0$ dla $n\ge1$.
Ma więc promień $+\infty$, chociaż
$\min(R_f,R_g)=1$.
Podobnie $f-f$ jest szeregiem zerowym.
\end{przyklad}
W przypadku algebry Banacha możemy jako $T$
wziąć mnożenie, nawet gdy algebra nie jest
przemienna. Kolejność czynników we współczynnikach
iloczynu jest wtedy istotna. W rozdziale
\ref{ma2:ch:banach-algebras}, twierdzenie
\ref{ma2:thm:algebra-power-series}, rozpatrzono
odmienną sytuację: skalarne współczynniki i argument
będący elementem algebry Banacha.

\section{Przykłady rozwinięć}

\begin{przyklad}
\label{ma2:ex:power-log-arctan}
Z szeregu geometrycznego i podstawienia $z=-t^2$
dla $|t|<1$ wynika
\[
 \frac1{1+t^2}=\sum_{n\ge0}(-1)^nt^{2n}.
\]
Całkując na odcinku o końcach $0,x$,
gdzie $|x|<1$, dostajemy
\[
 \arctan x=\sum_{n\ge0}\frac{(-1)^nx^{2n+1}}{2n+1}.
\]
Wykorzystano pochodną arcus tangens z AM I
(rozdział 18, sekcja 5) i wartość $\arctan0=0$.
Dla $x=1$ szereg po prawej stronie jest
zbieżny z kryterium Leibniza
\ref{ma1:thm:leibniz-series}, ponieważ
$1/(2n+1)$ maleje do zera. Twierdzenie
\ref{ma2:thm:power-abel-boundary} i ciągłość
arcus tangens dają
\[
 \frac{\pi}{4}=\sum_{n\ge0}\frac{(-1)^n}{2n+1}.
\]
Dokładnie tak samo z przykładu
\ref{ma2:ex:termwise-geometric-integral}
po podstawieniu $r$ zastąpionego przez $-x$
lub przez bezpośrednie całkowanie szeregu
$\sum_{n\ge0}(-t)^n$ otrzymujemy
$\log(1+x)=\sum_{n\ge1}(-1)^{n-1}x^n/n$
dla $|x|<1$. Dla $x=1$ kryterium Leibniza
i twierdzenie Abela dają sumę $\log2$.
\end{przyklad}

Powyższe rozwinięcia są podstawą dalszego
badania zespolonych funkcji elementarnych,
funkcji analitycznych i szeregów Taylora.
